\section{결과}
\label{sec:results}

\subsection{시공간 데이터베이스 표현으로의 재구성 결과}

RQ1에서 먼저 확인한 것은 모델 성능이 아니라 데이터의 단위다. PUBG 공식
텔레메트리의 시간순 이벤트 로그는 이동 객체, 시간변화 영역, 시간변화 속성, 시공간
사건, 위험집합으로 나눌 수 있었다. 그림~\ref{fig:team-phase-st-example}은
이 재구성이 단순한 분류표가 아니라 실제 맵 위의 시간 순서와 공간 변화를 함께
남긴다는 점을 보인다. 한 팀의 이동 궤적은 안전구역 경계의 변화와 함께 읽혀야 하며,
탈락 지점은 마지막 좌표라기보다 그 팀이 어떤 페이즈의 위험집합에서 빠졌는지를
가리킨다.

예측 단위도 여기서 바뀐다. 최종 순위 예측에서는 하나의 매치 또는 팀이 하나의 행으로
압축된다. 여기서는 페이즈 \(k\)까지 생존한 팀만 다음 페이즈의 위험집합에 남고,
라벨은 다음 페이즈에서의 탈락 여부가 된다. 테스트 요약 기준으로 보류 집합은 약
3.38k개의 페이즈별 생존 팀 라벨을 포함했고, 다음 페이즈 탈락 비율은 약
31.5\%였다. 이 구성은 PUBG 텔레메트리를 경기 후 집계 행이 아니라 페이즈별 생존
과정을 보존하는 시공간 데이터베이스 표현으로 다룰 수 있음을 확인한다.

\subsection{표현 구성요소 비교}

표~\ref{tab:component-ranking}은 \(X_s\)를 feature block 단위로 나누어 학습한 결과다.
표의 행들은 같은 종류의 ablation이 아니다. 팀 생존 여력, 이동 이력, 안전구역 맥락
단독은 특정 구성요소만 남긴 조건이고, 이동 이력+안전구역 맥락은 두 block을 단순히
붙인 조건이다. 최신 상태 기준모델은 시간 이력을 제거한 압축 기준이다. 이 구분을
먼저 두어야 결과를 과하게 읽지 않을 수 있다. 표의 평균 손실은 9개 맵/플레이 모드
조건별 테스트 BCE를 평균한 값이며, 표준편차는 조건 간 변동을 나타낸다.

\begin{table*}[t]
\centering
\caption{표현 구성요소 비교. 값은 조건별 테스트 BCE의 평균과 표준편차이며, 낮을수록 좋다.}
\label{tab:component-ranking}
\small
\setlength{\tabcolsep}{4pt}
\begin{tabular}{p{0.24\linewidth}p{0.45\linewidth}rr}
\toprule
조건 & 비교 의미 & 평균 손실 & 표준편차 \\
\midrule
팀 생존 여력
& 평균 체력, 최소 체력, 생존 인원, 차량 이용 여부로 구성한 페이즈별 시간변화 속성
& \textbf{0.5942} & 0.0206 \\
이동 이력
& 팀 중심 좌표와 속도의 10초 이력을 보존한 이동 객체 표현
& 0.6178 & 0.0181 \\
최신 상태 기준모델
& 페이즈 종료 시점의 단일 최신 상태만 사용하는 압축 기준
& 0.6311 & 0.0125 \\
이동 이력+안전구역 맥락
& 이동 이력에 안전구역 상대 위치와 경계 거리 변수를 단순 결합
& 0.6320 & 0.0157 \\
전체 팀 상태 스냅샷
& 팀 상태 스냅샷을 넓게 사용한 단순 확장 조건
& 0.6362 & 0.0214 \\
안전구역 맥락 단독
& 안전구역 중심, 경계 거리, 블루존 내부 여부 중심의 제한 표현
& 0.6411 & 0.0155 \\
\bottomrule
\end{tabular}
\end{table*}

팀 생존 여력은 평균 0.5942로 가장 낮은 손실을 보였다. 이 결과는 놀라운 발견이라기보다
전투 생존 게임에서 예상 가능한 직접 상태 변수의 효과다. 평균 체력, 최소 체력,
생존 인원, 차량 이용 여부는 그 팀이 다음 페이즈를 버틸 수 있는지를 거의 직접
나타낸다. 따라서 이 조건의 성능을 본 논문의 주된 기여로 해석하지 않는다.

시공간 표현 관점에서 더 중요한 비교는 이동 이력 조건이다. 이 조건은 체력, 생존
인원, 차량 이용 여부, 안전구역 변수를 제외하고 팀 중심 좌표와
속도의 10초 이력만 사용한다. 그럼에도 최신 상태 기준모델보다 낮은 손실을 보였다
(0.6178 vs. 0.6311). 이는 페이즈 종료 시점의 마지막 상태 하나만 남기는 것보다,
그 팀이 어떤 경로로 현재 위치에 도달했는지를 보존하는 편이 다음 페이즈 탈락 위험
예측에서 낫다는 뜻이다. 이 비교는 같은 페이즈 안의 사후 정보를 이용해 해당
페이즈의 생존을 맞힌 것이 아니라, 페이즈 \(k\) 종료 이전의 이력으로 페이즈
\(k+1\)의 탈락 여부를 예측한 결과다.

반면 안전구역 맥락 단독 조건은 가장 약했고, 이동 이력에 안전구역 변수를 단순
결합한 조건도 이동 이력만 사용한 조건보다 개선되지 않았다. 그래서 RQ2의 답은
부분적이다. 페이즈별 시간 정렬과 이동 이력 보존은 도움이 되었지만, 안전구역은
단독 표현이나 단순 결합만으로는 충분히 쓰이지 않았다.

\subsection{조건별 최저 손실과 후속 표현 설계}

표~\ref{tab:domain-winners}는 9개 맵/플레이 모드 조건에서 가장 낮은 테스트 손실을
보인 구성요소를 요약한다. 모든 조건에서 최저 손실을 보인 것은 팀 생존 여력이었다.
초기 비교 실험에서 먼저 잡힌 신호는 이동 객체의 좌표가 아니라, 페이즈별로 갱신되는
팀의 시간변화 속성이었다.

\begin{table}[t]
\centering
\caption{조건별 최저 테스트 손실. 모든 조건에서 팀 생존 여력이 최저 손실을 보였다.}
\label{tab:domain-winners}
\small
\begin{tabular}{llr}
\toprule
맵 & 플레이 모드 & 최저 손실 \\
\midrule
Baltic & TPP squad & 0.6038 \\
Tiger & TPP squad & 0.5608 \\
Desert & TPP squad & 0.5770 \\
Neon & TPP squad & 0.5788 \\
Savage & TPP squad & 0.5826 \\
Baltic & FPP squad & 0.5874 \\
Tiger & FPP squad & 0.6236 \\
Desert & FPP squad & 0.6148 \\
Neon & FPP squad & 0.6187 \\
\bottomrule
\end{tabular}
\end{table}

이 결과로 후속 표현 설계의 우선순위도 조금 분명해진다. 팀 생존 여력은 강하지만,
그 강함은 PUBG 스쿼드 매치가 전투 생존 게임이라는 사실에서 자연스럽게 나온다.
따라서 후속 연구의 질문은 생존 상태 변수가 중요하다는 것을 다시 확인하는 데 있지
않다. 더 중요한 것은 직접 생존 상태를 제외했을 때 이동 이력과 안전구역 맥락이
어떤 추가 정보를 제공하는지 분리하는 것이다. 이번 실험에서 이동
이력은 최신 상태 기준모델보다 나았으므로, 팀을 이동 객체로 보고 궤적을 보존하는
표현은 유지할 만하다. 반대로 안전구역은 상대 위치와 경계 거리 변수를 단독 또는
단순 결합으로 넣는 것만으로는 충분하지 않았다.

결과의 핵심은 특정 특징군의 승패가 아니다. PUBG 텔레메트리를 시공간 데이터베이스로
재구성하면 페이즈별 생존 예측 문제가 정의되고, 그 안에서 시간변화 속성과 이동
이력은 먼저 안정적인 신호로 나타난다. 특히 이동 이력만으로도 최신 상태 기준보다
낮은 손실을 보였다는 점은, 마지막 순간의 상태 하나보다 그 상태에 도달한 이력을
보존하는 것이 중요하다는 해석을 뒷받침한다. 안전구역은 다음 단계에서 이동 이력과
팀 생존 상태를 조건화하는 상호작용 표현으로 다시 설계해야 할 대상이다.
